\documentclass{article}
\usepackage[T1]{fontenc}
\usepackage{lmodern}
\usepackage{graphicx}
\usepackage{amsmath,amssymb}
\usepackage[a4paper,margin=2cm]{geometry}
\usepackage[absolute,overlay]{textpos}
\setlength{\TPHorizModule}{1mm}
\setlength{\TPVertModule}{1mm}
\usepackage[indonesian]{babel}
\usepackage{hyperref}
\setlength{\emergencystretch}{2em}
% unit: Halaman 85

\begin{document}

% === Halaman 85 (python/native) ===

• Contoh: 

    - misalkan bigram yang sering muncul di dalam cipherteks adalah BM. 

    - bigram yang sering muncul di dalam bahasa Inggris adalah TH

    - jadi, kita duga TH dipetakan menjadi BM

   - misalkan TH dienkripsi dengan aturan segiempat, maka jika T dan H membentuk 

segiempat dengan B dan M, maka posisi mereka harus konsisten dalam matriks.

     Artinya: T satu baris dengan B, dan H satu baris dengan M

                           T -------- B

                           |              |

                           M-------- H  

   - Jika TH dienkripsi dengan aturan kolom, maka B di bawah T, M di bawah H

   - Jika TH dienkripsi dengan aturan baris, maka B di kanan T, M di kanan H 

   - Mulai bangun matriks parsial 5 x 5 dengan hipotesis seperti di atas

85

\end{document}
